\section{Correctness of Context-Sharded Block Parallelism}
\label{app:correctness}

\begin{proposition}[Exactness of CSBP]
\label{prop:csbp-exactness}
Let $\{\mathcal{B}_r\}_{r=1}^{P}$ partition the $B$ target blocks, and let the
clean token rows be partitioned arbitrarily across the same $P$ ranks. Given
any fixed input sequence, sampled diffusion times and corruptions, per-token
model randomness, and loss normalization, CSBP computes, for every
parameter value $\theta$, the same global sampled loss and parameter gradient
as the reference length-$2L$ BDLM computation, up to floating-point reduction
order.
\end{proposition}

\begin{proof}
Fix the input sequence, sampled diffusion times and corruptions, model
randomness, and loss normalization. Let
$\mathcal L_b^{\mathrm{ref}}(\theta)$ and
$\mathcal L_b^{\mathrm{CSBP}}(\theta)$ be the resulting sampled losses for
block $b$.
For rank $r$, define
\begin{equation}
    \mathcal L_r^{\mathrm{CSBP}}
    =\sum_{b\in\mathcal B_r}\mathcal L_b^{\mathrm{CSBP}},
    \qquad
    \mathcal L^{\mathrm{CSBP}}
    =\sum_{r=1}^{P}\mathcal L_r^{\mathrm{CSBP}},
    \qquad
    \mathcal L^{\mathrm{ref}}
    =\sum_{b=1}^{B}\mathcal L_b^{\mathrm{ref}}.
    \label{eq:correctness-loss-definitions}
\end{equation}

Let $H_d^{\mathrm{clean}}$ and $H_{b,d}^{\mathrm{corr}}$ denote the clean
representations and the corrupted representations of block $b$ after $d$ of
the model's $D$ transformer layers. Thus $H_0$ is the embedding output and
$H_D$ enters the output projection. We use $H_d$ to denote all clean and
corrupted representations at depth $d$, and $\theta_d$ to denote the
parameters of layer $d$ for $1\leq d\leq D$.

\begin{enumerate}
\item
\textbf{Reference dependencies.}
The length-$2L$ construction and block-causal mask described in
Section~\ref{sec:background} determine the visible K/V rows. At layer $d$, a
clean query at position $q$ attends to clean positions $1,\ldots,q$. The
corrupted queries in block $b$ attend to their own corrupted block and the
clean positions in preceding blocks. We write $K_{\leq q,d}^{\mathrm{clean}}$
and $V_{\leq q,d}^{\mathrm{clean}}$ for the clean K/V rows at positions
$1,\ldots,q$. The reference attention is
\begin{align}
    O_{q,d}^{\mathrm{clean,ref}}
    &= \operatorname{Attn}\!\left(
        Q_{q,d}^{\mathrm{clean}},
        K_{\leq q,d}^{\mathrm{clean}},
        V_{\leq q,d}^{\mathrm{clean}}
      \right),
    \label{eq:correctness-clean-attention}
    \\
    O_{b,d}^{\mathrm{corr,ref}}
    &= \operatorname{Attn}\!\left(
        Q_{b,d}^{\mathrm{corr}},
        [K_{b,d}^{\mathrm{corr}};K_{<b,d}^{\mathrm{clean}}],
        [V_{b,d}^{\mathrm{corr}};V_{<b,d}^{\mathrm{clean}}]
      \right).
    \label{eq:correctness-corrupted-attention}
\end{align}
Equation~\ref{eq:correctness-corrupted-attention} implies, for every
$Z\in\{Q,K,V\}$,
\begin{equation}
    \frac{\partial O_{b,d}^{\mathrm{corr,ref}}}
         {\partial Z_{b',d}^{\mathrm{corr}}}
    =0
    \qquad (b'\neq b).
    \label{eq:correctness-block-dependency}
\end{equation}

\item
\textbf{Exact sharded attention.}
Let $C_r$ be the set of clean token positions stored on rank $r$. The clean
shards satisfy
\begin{equation}
    \bigcup_{r=1}^{P} C_r=\{1,\ldots,L\},
    \qquad
    C_r\cap C_{r'}=\varnothing \quad (r\neq r').
    \label{eq:correctness-clean-partition}
\end{equation}
Consider one query vector $\mathbf q$. Let $S$ be its visible key positions,
partitioned into disjoint nonempty sets $S_1,\ldots,S_J$. For the key and value
rows $\mathbf k_i$ and $\mathbf v_i$, let
$a_i=\mathbf q\mathbf k_i^\top/\sqrt{d_h}$ and
$a_{\max}=\max_{i\in S}a_i$. Since the sets $S_j$ partition $S$,
\begin{equation}
    \operatorname{Attn}(\mathbf q,K_S,V_S)
    =
    \frac{\displaystyle\sum_{j=1}^{J}\sum_{i\in S_j}
          e^{a_i-a_{\max}}\mathbf v_i}
         {\displaystyle\sum_{j=1}^{J}\sum_{i\in S_j}e^{a_i-a_{\max}}}.
    \label{eq:correctness-online-softmax-merge}
\end{equation}
Online softmax evaluates these partitioned sums exactly
\citep{milakov2018online,dao2022flashattention}. For a clean query at position
$q$, the nonempty sets $C_r\cap\{1,\ldots,q\}$ partition its visible K/V rows.
For a corrupted query in block $b$, its local corrupted block together with
the nonempty clean sets $C_r\cap\{1,\ldots,(b-1)M\}$ partition its visible
K/V rows. These clean sets are empty for $b=1$.
CSBP therefore supplies exactly the rows in
Equations~\ref{eq:correctness-clean-attention} and
\ref{eq:correctness-corrupted-attention}.

\item
\textbf{Forward induction.}
The common inputs and embedding computation give
$H_0^{\mathrm{CSBP}}=H_0^{\mathrm{ref}}$. Assume
$H_{d-1}^{\mathrm{CSBP}}=H_{d-1}^{\mathrm{ref}}$ for some
$1\leq d\leq D$. Their Q/K/V projections agree, and
Equation~\ref{eq:correctness-online-softmax-merge} gives
\begin{equation}
    \begin{aligned}
      O_{q,d}^{\mathrm{clean,CSBP}}&=O_{q,d}^{\mathrm{clean,ref}}
      && \forall q,\\
      O_{b,d}^{\mathrm{corr,CSBP}}=O_{b,d}^{\mathrm{corr,ref}}
      && \forall b.
    \end{aligned}
    \label{eq:correctness-forward-step}
\end{equation}
The remaining operations in layer $d$ are unchanged, so
$H_d^{\mathrm{CSBP}}=H_d^{\mathrm{ref}}$. Induction over
$d=1,\ldots,D$ gives $H_D^{\mathrm{CSBP}}=H_D^{\mathrm{ref}}$.

\item
\textbf{Loss.}
The output projection and block losses are unchanged. The forward induction
therefore gives
$\mathcal L_b^{\mathrm{CSBP}}=\mathcal L_b^{\mathrm{ref}}$ for every block.
Since $\{\mathcal B_r\}_{r=1}^{P}$ partitions the target blocks,
\begin{equation}
    \mathcal L^{\mathrm{CSBP}}
    =\sum_{r=1}^{P}\sum_{b\in\mathcal B_r}\mathcal L_b^{\mathrm{CSBP}}
    =\sum_{b=1}^{B}\mathcal L_b^{\mathrm{ref}}
    =\mathcal L^{\mathrm{ref}}.
    \label{eq:correctness-sampled-loss}
\end{equation}

\item
\textbf{Backward pass.}
The equal final representations and common loss functions give the same
derivatives at the output of the model. Let $\mathcal L$ denote the global
sampled loss of either execution. For
$Z\in\{Q,K,V\}$,
Equation~\ref{eq:correctness-block-dependency} and the chain rule give
\begin{equation}
    \frac{\partial \mathcal L}{\partial Z_{b,d}^{\mathrm{corr}}}
    = \sum_{b'=1}^{B}
      \left(
      \frac{\partial O_{b',d}^{\mathrm{corr}}}
           {\partial Z_{b,d}^{\mathrm{corr}}}
      \right)^{\!\top}
      \frac{\partial \mathcal L}{\partial O_{b',d}^{\mathrm{corr}}}
    = \left(
      \frac{\partial O_{b,d}^{\mathrm{corr}}}
           {\partial Z_{b,d}^{\mathrm{corr}}}
      \right)^{\!\top}
      \frac{\partial \mathcal L}{\partial O_{b,d}^{\mathrm{corr}}}.
    \label{eq:correctness-corrupted-gradient}
\end{equation}
Thus corrupted Q/K/V derivatives depend only on the queries from the same
block and are computed on its owner. Let $Z_{r,d}^{\mathrm{clean}}$ denote the
clean K/V rows at positions $C_r$, and let $O_{q,d}$ denote the attention
output for query row $q$. The chain rule gives
\begin{equation}
    \frac{\partial \mathcal L}{\partial Z_{r,d}^{\mathrm{clean}}}
    = \sum_q
      \left(
        \frac{\partial O_{q,d}}{\partial Z_{r,d}^{\mathrm{clean}}}
      \right)^{\!\top}
      \frac{\partial \mathcal L}{\partial O_{q,d}},
    \qquad Z\in\{K,V\},
    \label{eq:correctness-clean-gradient}
\end{equation}
where the sum is over all clean and corrupted query rows. A term is zero when
the query cannot attend to a position in $C_r$. Exact CP backward accumulates
this sum on rank $r$. For every clean query row $q$,
\begin{equation}
    \frac{\partial \mathcal L}{\partial Q_{q,d}^{\mathrm{clean}}}
    = \left(
        \frac{\partial O_{q,d}^{\mathrm{clean}}}
             {\partial Q_{q,d}^{\mathrm{clean}}}
      \right)^{\!\top}
      \frac{\partial \mathcal L}{\partial O_{q,d}^{\mathrm{clean}}},
    \label{eq:correctness-query-gradient}
\end{equation}
which is computed on the query owner. Because the two executions have the same
forward values and output derivatives, Equations~\ref{eq:correctness-corrupted-gradient}--\ref{eq:correctness-query-gradient}
give the same Q/K/V derivatives. Backpropagation through the unchanged
remainder of layer $d$ therefore satisfies
\begin{equation}
    \begin{aligned}
      \frac{\partial \mathcal L^{\mathrm{CSBP}}}
           {\partial H_{d-1}^{\mathrm{CSBP}}}
      &=
      \left(
        \frac{\partial H_d^{\mathrm{CSBP}}}
             {\partial H_{d-1}^{\mathrm{CSBP}}}
      \right)^{\!\top}
      \frac{\partial \mathcal L^{\mathrm{CSBP}}}
           {\partial H_d^{\mathrm{CSBP}}} \\
      &=
      \left(
        \frac{\partial H_d^{\mathrm{ref}}}
             {\partial H_{d-1}^{\mathrm{ref}}}
      \right)^{\!\top}
      \frac{\partial \mathcal L^{\mathrm{ref}}}
           {\partial H_d^{\mathrm{ref}}}
      =
      \frac{\partial \mathcal L^{\mathrm{ref}}}
           {\partial H_{d-1}^{\mathrm{ref}}}, \\[4pt]
      \nabla_{\theta_d}\mathcal L^{\mathrm{CSBP}}
      &=
      \left(
        \frac{\partial H_d^{\mathrm{CSBP}}}{\partial \theta_d}
      \right)^{\!\top}
      \frac{\partial \mathcal L^{\mathrm{CSBP}}}
           {\partial H_d^{\mathrm{CSBP}}} \\
      &=
      \left(
        \frac{\partial H_d^{\mathrm{ref}}}{\partial \theta_d}
      \right)^{\!\top}
      \frac{\partial \mathcal L^{\mathrm{ref}}}
           {\partial H_d^{\mathrm{ref}}}
      =\nabla_{\theta_d}\mathcal L^{\mathrm{ref}}.
    \end{aligned}
    \label{eq:correctness-backward-step}
\end{equation}
Because the output projection is identical, its parameter gradient and
derivative with respect to $H_D$ agree, providing the base case. Reverse
induction through $d=D,\ldots,1$ gives equal transformer-layer gradients and
equal derivatives with respect to $H_0$. The common embedding computation then
gives equal embedding gradients. Finally, linearity of differentiation and
Equation~\ref{eq:correctness-loss-definitions} give
\begin{equation}
    \sum_{r=1}^{P}
    \nabla_\theta \mathcal L_r^{\mathrm{CSBP}}
    =\nabla_\theta \mathcal L^{\mathrm{CSBP}}
    =\nabla_\theta \mathcal L^{\mathrm{ref}}.
    \label{eq:correctness-sampled-gradient}
\end{equation}
\end{enumerate}
Equations~\ref{eq:correctness-sampled-loss} and
\ref{eq:correctness-sampled-gradient} hold for every sampled realization.
Taking expectations also gives the same BDLM objective and objective gradient.
\end{proof}

\paragraph{Communication consequence.}
Under CSBP, only clean K/V and their gradients cross ranks during attention,
while corrupted Q/K/V and their gradients remain on the rank assigned their
block.

\section{Load Balancing}
\label{app:load-balancing}

Both attention workloads in CSBP increase with block position. Later
target blocks attend to longer clean prefixes, and later clean query blocks
attend to longer causal prefixes. Because both costs grow with the same block
index, the same early--late pairing balances them simultaneously.

Assume $B$ equal-sized blocks of $M=L/B$ tokens. Let
$W_b^{\mathrm{target}}$ and $W_b^{\mathrm{clean}}$ denote the valid query--key
pairs for target-block and clean causal attention at block $b$. The $M$ target
queries in block $b$ attend to $(b-1)M$ clean-prefix keys and $M$ keys from
their own corrupted block, for $bM$ visible keys. The $M$ clean queries attend
to the $b-1$ preceding clean blocks and causally within their own block. Their
work is therefore
\begin{align}
    W_b^{\mathrm{target}}
    &= bM^2,
    &
    W_b^{\mathrm{clean}}
    &= (b-1)M^2+\frac{M(M+1)}{2}.
    \label{eq:block-work-by-position}
\end{align}

Dual-end BP pairs target block $b$ with block $B+1-b$. Clean-sequence zigzag
sharding uses the same early--late pairing~\citep{dubey2024llama3}. The work of
each pair is
\begin{align}
    W_b^{\mathrm{target}}+W_{B+1-b}^{\mathrm{target}}
    &= (B+1)M^2,
    &
    W_b^{\mathrm{clean}}+W_{B+1-b}^{\mathrm{clean}}
    &= BM^2+M.
    \label{eq:aligned-pair-work}
\end{align}
Both sums are independent of $b$, so every early--late pair carries the same
amount of target-block work and the same amount of clean causal work. When $B$
is divisible by $2P$, assigning $B/(2P)$ pairs to each rank gives
\begin{align}
    W_r^{\mathrm{target}}
    &= \frac{B(B+1)M^2}{2P},
    &
    W_r^{\mathrm{clean}}
    &= \frac{L(L+1)}{2P},
    \qquad 1\leq r\leq P.
    \label{eq:aligned-rank-work}
\end{align}
Using the same pairs for target-block placement and clean zigzag sharding gives
every rank equal target-block and clean causal attention work. Attention
backward uses the same valid query--key pairs, so these pair counts remain
balanced. When exact divisibility does not hold, we distribute the complete
pairs as evenly as possible and assign any unpaired middle block to the
currently least-loaded rank.

\clearpage
\section{CSBP Algorithm and Execution Schedule}
\label{app:execution-schedule}

Algorithm~\ref{alg:csbp-attention} summarizes one CSBP layer. Each rank stores
one clean-sequence shard and the complete corrupted blocks in $\mathcal B_r$.
All ranks execute concurrently: CP supplies clean K/V, while each block's
corrupted Q/K/V remain local.

\begin{algorithm}[H]
\small
\caption{CSBP forward pass on rank $r$}
\label{alg:csbp-attention}
\begin{algorithmic}[1]
\State \textbf{Input:} Clean shard $(Q_r,K_r,V_r)$ and
  $\{(Q_b,K_b,V_b):b\in\mathcal B_r\}$
\State \textbf{Output:} $O_r$ and $\{O_b:b\in\mathcal B_r\}$
\State All-gather the clean K/V shards to form
  $(K^{\mathrm{clean}},V^{\mathrm{clean}})$ on each rank
\Statex
\State \textbf{Clean-sequence attention (CP)}
\State Compute clean-query attention under the block-causal mask
\State $O_r\gets\operatorname{softmax}\!\left(
  Q_r(K^{\mathrm{clean}})^\top/\sqrt{d_h}+\mathcal M_r\right)
  V^{\mathrm{clean}}$
\Statex
\State \textbf{Assigned corrupted blocks (BP)}
\ForAll{$b\in\mathcal B_r$}
  \State Compute within-block attention from the local $Q_b,K_b,V_b$
  \State Compute clean-prefix attention from the gathered
    $K_{<b}^{\mathrm{clean}},V_{<b}^{\mathrm{clean}}$
  \State Merge the two outputs and LSEs by online softmax to obtain
  \State $O_b\gets\operatorname{Attn}\!\left(
    Q_b,[K_b;K_{<b}^{\mathrm{clean}}],
    [V_b;V_{<b}^{\mathrm{clean}}]\right)$
    \Comment{Equation~\ref{eq:csbp-block-attention}}
\EndFor
\State \Return $O_r$ and $\{O_b:b\in\mathcal B_r\}$
\end{algorithmic}
\end{algorithm}

Algorithm~\ref{alg:csbp-backward} gives the exact backward pass. Rank $r$
keeps the corrupted Q/K/V gradients for its assigned blocks, while CP sums the
clean K/V gradients contributed by all clean and corrupted queries.

\begin{algorithm}[H]
\small
\caption{CSBP attention backward on rank $r$}
\label{alg:csbp-backward}
\begin{algorithmic}[1]
\State \textbf{Input:} Saved forward tensors and
  $\{\mathcal L_b:b\in\mathcal B_r\}$
\State \textbf{Output:} $(dQ_r,dK_r,dV_r)$ and
  $\{(dQ_b,dK_b,dV_b):b\in\mathcal B_r\}$
\Statex
\State \textbf{Clean-query backward (CP)}
\State Backpropagate $O_r$; retain $dQ_r$ and the clean K/V contributions
\Statex
\State \textbf{Assigned corrupted blocks (local backward)}
\ForAll{$b\in\mathcal B_r$}
  \State Backpropagate $\mathcal L_b$ through $O_b$
  \State $(dQ_b,dK_b,dV_b)
    \gets\nabla_{(Q_b,K_b,V_b)}\mathcal L_b$
    \Comment{kept on rank $r$}
  \State Add this block's clean-prefix K/V gradients to the CP reduction
\EndFor
\Statex
\State \textbf{Clean-gradient reduction (CP backward)}
\State Reduce-scatter the clean K/V contributions across ranks
\For{$Z\in\{K,V\}$}
  \State $\displaystyle dZ_r\gets
    \sum_q\left(\frac{\partial O_q}{\partial Z_r}\right)^{\!\top}
    \frac{\partial\mathcal L}{\partial O_q}$
    \Comment{Equation~\ref{eq:correctness-clean-gradient}}
\EndFor
\State \Return $(dQ_r,dK_r,dV_r)$ and
  $\{(dQ_b,dK_b,dV_b):b\in\mathcal B_r\}$
\end{algorithmic}
\end{algorithm}

\clearpage
\paragraph{Efficient fused operator.}
Implementation~\ref{impl:csbp-forward} shows the production forward path. It
starts an asynchronous all-gather of clean K/V, then computes all local
corrupted blocks with one block-diagonal attention launch while communication
runs. After the all-gather, one masked fused-attention launch computes the
clean-K/V contribution for every query. The mask uses each gathered key's
original clean token position, so the kernel consumes the gathered K/V directly
without a reordering copy.

Each attention launch returns a normalized output $O_j$ and row-wise
log-sum-exp (LSE) $z_j$.  For two disjoint key sets, the CUDA kernel sets
$m=\max(z_1,z_2)$, $w_j=e^{z_j-m}$, and computes
\begin{equation}
  O=\frac{w_1O_1+w_2O_2}{w_1+w_2},\qquad
  z=m+\log(w_1+w_2).
  \label{eq:csbp-lse-merge}
\end{equation}
This is Equation~\ref{eq:correctness-online-softmax-merge} applied to the two
attention launches.  The implementation stores the numerator, $m$, and the
normalizer $w_1+w_2$ in FP32 until the final division.

\begin{lstlisting}[
  style=implementation,
  language=Python,
  caption={CSBP forward schedule and CUDA LSE merge.},
  label={impl:csbp-forward}
]
@cuda_kernel
def merge_lse_kernel(numerator, m, normalizer,
                     O_block, z_block):
    # One thread handles one (corrupted-query row, head-dimension) element.
    q, d = cuda_grid_2d()
    if q >= len(z_block):
        return
    m_new = max(m[q], z_block[q])
    w_acc = exp(m[q] - m_new) if isfinite(m[q]) else 0.0
    w_block = exp(z_block[q] - m_new) if isfinite(z_block[q]) else 0.0
    numerator[q, d] = (w_acc * numerator[q, d]
                       + w_block * O_block[q, d])
    if d == 0:
        normalizer[q] = w_acc * normalizer[q] + w_block
        m[q] = m_new


def csbp_collective_forward(Q, K_block, V_block,
                            K_clean_r, V_clean_r, masks, scale):
    # 1. Start the clean K/V all-gather.
    work, gathered = async_all_gather(K_clean_r, V_clean_r)

    # 2. One block-diagonal launch covers every block owned by this rank.
    O_block, LSE_block = block_diagonal_attention(
        Q[:masks.corrupted_rows], K_block, V_block, scale)

    # 3. Use the gathered clean K/V directly in one masked launch.
    work.wait()
    K_clean, V_clean = view_all_gathered_clean_kv(gathered)
    clean_mask = block_causal_mask(masks)
    O_clean, LSE_clean = masked_clean_attention(
        Q, K_clean, V_clean, clean_mask, scale)

    # 4. Custom CUDA kernels apply the LSE rule above in FP32.
    numerator = O_clean.float()
    m = LSE_clean.float()
    normalizer = isfinite(m).float()
    launch(merge_lse_kernel, numerator, m, normalizer,
           O_block, LSE_block)
    O, LSE = empty_outputs(Q)
    launch(normalize_output_kernel,
           numerator, m, normalizer, O, LSE)
    save_for_backward(Q, K_block, V_block, K_clean, V_clean, O, LSE)
    return O
\end{lstlisting}

\clearpage
Implementation~\ref{impl:csbp-backward} shows the fused packed backward. One
masked kernel reuses the saved output and LSE to compute gradients for the local
corrupted K/V and gathered clean K/V together. The corrupted gradients remain
local, while only the clean K/V gradients are reduce-scattered to their CP
ranks.

\begin{lstlisting}[
  style=implementation,
  language=Python,
  caption={Fused CSBP backward. One masked kernel computes the local corrupted-block and clean-context gradients.},
  label={impl:csbp-backward}
]
def csbp_fused_backward(dO, saved, masks, scale):
    Q, K_block, V_block, K_clean, V_clean, O, LSE = saved

    # 1. Pack local corrupted K/V before the gathered clean K/V.
    K_packed, V_packed = pack_local_clean_kv(
        K_block, V_block, K_clean, V_clean)

    # 2. One masked backward reuses the exact forward output and LSE.
    dQ, dK_packed, dV_packed = masked_attention_backward(
        Q, K_packed, V_packed, O, LSE, dO,
        masks.packed, scale)

    # 3. Split local gradients from the clean gradients that communicate.
    dK_block, dK_clean = split_local_clean_gradients(
        dK_packed, masks.corrupted_rows)
    dV_block, dV_clean = split_local_clean_gradients(
        dV_packed, masks.corrupted_rows)

    # 4. Return each rank's clean K/V-gradient shard.
    dK_clean_r, dV_clean_r = reduce_scatter_clean_gradients(
        dK_clean, dV_clean)
    return dQ, dK_block, dV_block, dK_clean_r, dV_clean_r
\end{lstlisting}

\section{AR-to-BDLM Conversion Methodology}
\label{app:fast-dllm-v2-objective}

Fast-dLLM v2 adapts a pretrained causal LM to block diffusion while retaining
next-token prediction~\citep{wu2025fastdllmv2}. It samples a blockwise binary
mask $m^{(0)}$ and constructs a second noisy view with the complementary mask
$m^{(1)}=1-m^{(0)}$ on every supervised token. If $h_i^{(v)}$ is the hidden
representation at position $i$ in view $v\in\{0,1\}$, let $m_i^{(v)}=1$ when
$x_i$ is masked and let $\mathcal{S}$ contain the supervised positions other
than the first token. The sampled objective is
\begin{equation}
    \widehat{\mathcal{L}}_{\mathrm{AR\mbox{-}to\mbox{-}BDLM}}
    =
    -\sum_{v=0}^{1}\sum_{i\in\mathcal{S}}
    m_i^{(v)}\log p_\theta\!\left(x_i\mid h_{i-1}^{(v)}\right),
    \qquad
    m_i^{(0)}+m_i^{(1)}=1.
    \label{eq:fast-dllm-v2-objective}
\end{equation}
The complementary views supervise every eligible next token exactly once.
The one-token shift uses the preceding query row to predict $x_i$, matching a
causal LM head.

The objective preserves the two properties used by block parallelism. Each
view uses the length-$2L$ clean-plus-corrupted computation, and queries in
block $b$ use corrupted K/V only from block $b$ of that view and clean K/V
from preceding clean blocks,
\begin{equation}
    O_{b,v}
    =\operatorname{Attn}\!\left(
      Q_{b,v},
      [K_{b,v};K_{<b,v}^{\mathrm{clean}}],
      [V_{b,v};V_{<b,v}^{\mathrm{clean}}]
    \right).
    \label{eq:fast-dllm-v2-attention}
\end{equation}
The two views are separate batch elements and introduce no cross-view
attention. The label shift also leaves these attention dependencies unchanged.
We assign each shifted loss term to the block containing its prediction row,
and group the terms from both views into the corresponding block loss. The
objective therefore retains the decomposition
$\widehat{\mathcal{L}}_{\mathrm{AR\mbox{-}to\mbox{-}BDLM}}
=\sum_{b=1}^{B}\mathcal{L}_b$.

CSBP therefore applies directly. BP assigns each complete corrupted
computation and loss for block $b$ in both views to one rank, while CP shards
the clean sequence of both views across the same ranks.
Corrupted Q/K/V and their gradients remain on the block owner, while clean K/V
gradients are accumulated on their CP shard owners.
Proposition~\ref{prop:csbp-exactness}
applies to each view, and summing the complementary-view losses preserves its
loss and gradient equalities.

\clearpage
\section{DFlash2 Drafter Training Methodology}
\label{app:dflash2-objective}

DFlash trains a lightweight block-diffusion drafter for speculative
decoding~\citep{chen2026dflash}. Given a clean sequence $x$, hidden features
$H^{\mathrm{tar}}$ are extracted from a frozen autoregressive target model.
For each sampled block $b$, the token at position $a_b$ is visible and the next
$M-1$ tokens are masked and predicted in one pass. Queries attend
bidirectionally within the block and to target-feature K/V preceding $a_b$;
different draft blocks do not attend to one another.

Let
$\widetilde{x}^{(b)}=(x_{a_b},\mathtt{[MASK]},\ldots,\mathtt{[MASK]})$
denote the input for block $b$. One draft-attention layer computes
\begin{equation}
    O_b
    =
    \operatorname{Attn}\!\left(
      Q_b,
      [K_b;K_{<a_b}^{\mathrm{tar}}],
      [V_b;V_{<a_b}^{\mathrm{tar}}]
    \right).
    \label{eq:dflash2-block-attention}
\end{equation}
Following DFlash2~\citep{inco2026dflash2}, our drafter uses the architecture in
the open-source SpecForge implementation~\citep{specforge2025}. It retains this
one-pass backbone and uses block-local two-tap dynamic convolutions and a
candidate-path selector. Its sampled objective is the sum of a draft-token loss
and a selector loss for each block:
\begin{equation}
    \widehat{\mathcal{L}}_{\mathrm{DFlash2}}
    = \frac{1}{Z}\sum_b \mathcal{L}_b,
    \qquad
    \mathcal{L}_b
    = \mathcal{L}^{\mathrm{draft}}_b
      + \alpha_{\mathrm{sel}}\mathcal{L}^{\mathrm{sel}}_b,
    \label{eq:dflash2-training-objective}
\end{equation}
where $Z$ is the sum of valid token weights and $\alpha_{\mathrm{sel}}$ weights
the selector term.

This blockwise objective makes CSBP directly applicable. Let
$\mathcal B_r$ be the blocks assigned to rank $r$; these sets are disjoint and
together contain every sampled block. Since $\mathcal L_b$ depends only on block
$b$ and its target-feature prefix, each rank evaluates its assigned terms, and
\begin{equation}
    \sum_{r=1}^{P}
    \nabla_\theta\!\left(
      \frac{1}{Z}\sum_{b\in\mathcal B_r}\mathcal L_b
    \right)
    = \nabla_\theta\widehat{\mathcal L}_{\mathrm{DFlash2}}.
    \label{eq:dflash2-gradient-decomposition}
\end{equation}
Thus, assigning complete blocks to ranks preserves the loss and gradient up to
floating-point reduction order. It also gives DFlash2 an additional source of
savings beyond full-model SFT and AR-to-BDLM conversion. With $P$ context
ranks, conventional CP repeats each block-local drafter path $P$ times, while
CSBP evaluates it once. Full-model training has no separate drafter path and
therefore no corresponding $P$-to-one reduction. The repeated DFlash2 work
includes the decoder, vocabulary objective, output-head recomputation, and
selector; target-feature processing remains context sharded. At 512K, removing
two-way replication reduces measured GEMM work by $1.92\times$, and higher
execution efficiency raises the throughput gain to $2.48\times$. At 1M,
removing eight-way replication permits four data-parallel replicas, producing
the $7.59\times$ aggregate throughput gain.

\section{DiffusionGemma SFT Training Data and Evaluation Details}
\label{app:training-evaluation-details}

\paragraph{SWE-bench Verified.}
For the software-engineering comparison in
Figure~\ref{fig:equal-wall-clock-downstream}, the best baseline and CSBP
train on identical
ordered streams of successful CoderForge-Preview trajectories
\citep{CoderForge2026}, formatted for the OpenHands agent
interface~\citep{wang2025openhands}. Trajectories in the training stream have
a mean sequence length of 65{,}536.5 tokens and a maximum length of
131{,}072 tokens.
We evaluate the base model at 0 hours and
checkpoints at 3, 6, 9, and 12 hours, with one attempt on each of the 500
SWE-bench Verified tasks~\citep{jimenez2024swebench,chowdhury2024swebenchverified}.
The harness supplies the issue and repository state; success requires both the
designated \texttt{FAIL\_TO\_PASS} and \texttt{PASS\_TO\_PASS} tests to pass.

\paragraph{Terminal-Bench Lite.}
The best baseline and CSBP train on identical ordered streams of quality-filtered Terminus-2
trajectories from
LiteCoder-Terminal-SFT~\citep{peng2026litecoderterminal}. Trajectories in the
training stream have a mean sequence length of 65{,}551 tokens and a maximum
length of 131{,}072 tokens. We evaluate the same
checkpoint schedule with Harbor's Terminus-2 agent~\citep{Harbor_Framework}
on all 100 OpenThoughts-TBLite tasks~\citep{OpenThoughts-TBLite}, with one
attempt per task. A task passes only when its benchmark verifier returns success.

\paragraph{Matched wall-clock protocol.}
Within each workload, the best baseline and CSBP use the same model, training corpus,
optimizer settings, checkpoint schedule, and evaluation configuration; only
the distributed training method changes. We report pass rate as the percentage
of tasks solved by the single attempt. Evaluation tasks and verifier outputs do
not contribute training gradients. Both sources apply benchmark decontamination:
CoderForge excludes matching SWE-bench Verified repository/commit pairs and
problem statements, while LiteCoder-Terminal-SFT filters 13-gram overlaps with
Terminal-Bench queries~\citep{CoderForge2026,peng2026litecoderterminal}.

\paragraph{SFT optimization.}
The CoderForge-Preview and LiteCoder-Terminal-SFT runs use a global batch size
of one example and one gradient-accumulation step per update.
We train rank-16 LoRA adapters with $\alpha=32$ on eight H100 GPUs. The
baseline uses CP8, while CSBP uses CP8/BP8.
We optimize with DeepSpeed ZeRO-2 and FusedAdam using a peak learning rate of
$10^{-5}$, Adam coefficients $(\beta_1,\beta_2)=(0.95,0.99)$,
$\epsilon=10^{-8}$, weight decay $10^{-4}$, gradient clipping at $1.0$, and a
cosine schedule that decays the learning rate to $10^{-6}$.

\section{End-to-End Numerical Agreement}
\label{app:numerical-agreement}

Table~\ref{tab:objective-equivalence-256k} compares the strongest non-BP
baseline with CSBP using full-depth BF16 runs at 256K context on sixteen H200
GPUs.  The
NemotronDiffusion and DiffusionGemma rows evaluate BDLM fine-tuning, while the
Qwen rows evaluate Fast-dLLM v2 AR-to-BDLM conversion.  Each pair uses the same
model, random seed, generated inputs, objective, and optimizer configuration.
The resulting loss differences are small and consistent with floating-point
reduction order.

% Losses are rank means averaged over the measured steps in the selected runs.
% NemotronDiffusion 3B (two measured steps):
% /persistent/runs/nemotron_labs_diffusion_3b_262144_block256_h200_nodes2_tp1_ep1_cp4_bp4_mb1_xnodecp_20260830_201529
% NemotronDiffusion 8B (two measured steps):
% /persistent/runs/nemotron_labs_diffusion_8b_262144_block256_h200_nodes2_tp1_ep1_cp4_bp4_mb1_xnodecp_20260813_162942
% NemotronDiffusion 14B (two measured steps):
% /persistent/runs/nemotron_labs_diffusion_14b_262144_block256_h200_nodes2_tp1_ep1_cp8_bp8_mb1_xnodecp_20260813_170220
% DiffusionGemma 26B-A4B (five measured steps):
% /persistent/runs/diffusiongemma_26b_a4b_it_262144_block256_h200_nodes2_tp1_ep2_cp8_bp8_mb1_scheduledual_end_xnodecp_20260914_001756
% Qwen3.8-27B Fast-dLLM v2 (two measured steps):
% /persistent/runs/qwen3.8_27b_262144_h200_fast_dllm_v2_nodes2_tp2_cp8_20260909_172222
% Qwen3.5-27B Fast-dLLM v2 (two measured steps; selected baseline / ours):
% /persistent/runs/qwen3.5_27b_262144_h200_fast_dllm_v2_nodes2_tp2_cp8_20260830_110221
% /persistent/runs/qwen3.5_27b_262144_h200_fast_dllm_v2_nodes2_tp2_cp8_20260830_113949
\begin{table}[H]
  \centering
  \caption{End-to-end loss agreement between the strongest non-BP baseline and CSBP at 256K context.}
  \label{tab:objective-equivalence-256k}
  \small
  \setlength{\tabcolsep}{5pt}
  \renewcommand{\arraystretch}{1.12}
  \resizebox{\linewidth}{!}{%
  \begin{tabular}{@{}lcllrr@{}}
    \toprule
    Model & Context & \shortstack{Best baseline\\topology} & \shortstack{Ours\\topology} & \shortstack{Absolute loss\\difference} & \shortstack{Relative loss\\difference} \\
    \midrule
    NemotronDiffusion 3B
      & 256K & DP4/TP1/CP4 & DP4/TP1/CP4/BP4
      & $2.83\times10^{-3}$ & $2.11\times10^{-4}$ \\
    NemotronDiffusion 8B
      & 256K & DP4/TP1/CP4 & DP4/TP1/CP4/BP4
      & $2.59\times10^{-3}$ & $1.97\times10^{-4}$ \\
    NemotronDiffusion 14B
      & 256K & DP2/TP1/CP8 & DP2/TP1/CP8/BP8
      & $9.54\times10^{-7}$ & $7.18\times10^{-8}$ \\
    DiffusionGemma 26B-A4B
      & 256K & DP1/TP1/EP2/CP8 & DP1/TP1/EP2/CP8/BP8
      & $6.51\times10^{-3}$ & $2.11\times10^{-4}$ \\
    Qwen3.8-27B
      & 256K & DP1/TP2/CP8 & DP1/TP2/CP8/BP8
      & $7.21\times10^{-6}$ & $5.11\times10^{-7}$ \\
    Qwen3.5-27B
      & 256K & DP1/TP2/CP8 & DP1/TP2/CP8/BP8
      & $9.50\times10^{-5}$ & $6.69\times10^{-6}$ \\
    \bottomrule
  \end{tabular}}
\end{table}


\section{Pure Block Parallelism Ablation}
\label{app:pure-bp-ablation}

Sharding the clean context makes BP practical at long sequence lengths
(Table~\ref{tab:pure-bp-ablation}). For NemotronDiffusion 14B at 64K, CSBP is
$1.46\times$ faster than pure BP and uses 89.8 rather than 110.7\,GiB of HBM.
At 128K, CSBP still fits while pure BP runs out of memory. DiffusionGemma 26B-A4B
shows the same need for context sharding even earlier: pure BP runs out of
memory at both 64K and 128K, where CSBP fits and outperforms pure CP. Block ownership removes corrupted
K/V communication; context sharding prevents replicated clean prefixes from
erasing that benefit.

% Pure-BP provenance:
% NemotronDiffusion 14B, 64K (completed):
% /persistent/runs/nemotron_labs_diffusion_14b_65536_block256_h200_nodes2_tp1_ep1_cp2_bp2_mb1_20260831_213357/bp
% NemotronDiffusion 14B, 128K (measured OOM in checkpointed MLP backward):
% /persistent/runs/nemotron_labs_diffusion_14b_131072_block256_h200_nodes2_tp1_ep1_cp4_bp4_mb1_20260901_113051/bp
% DiffusionGemma 26B-A4B native all-block objective, 64K (measured OOM in backward):
% /persistent/runs/diffusiongemma_26b_a4b_it_65536_block256_h200_nodes2_tp1_ep2_cp2_bp2_mb1_scheduledual_end_20260914_025123/bp
% DiffusionGemma 26B-A4B native all-block objective, 128K (measured OOM in backward):
% /persistent/runs/diffusiongemma_26b_a4b_it_131072_block256_h200_nodes2_tp1_ep2_cp4_bp4_mb1_scheduledual_end_20260914_025127/bp
\begin{table}[H]
  \centering
  \caption{Context sharding enables block parallelism at long contexts.}
  \label{tab:pure-bp-ablation}
  \footnotesize
  \setlength{\tabcolsep}{1.75pt}
  \renewcommand{\arraystretch}{1.08}
  % ICLR's larger footnote font otherwise makes this table exceed the column.
  % The NeurIPS branch is an identity, preserving its existing rendering.
  \ifdefined\iclrfinalcopy
    \def\fitpurebptable#1{\resizebox{\linewidth}{!}{#1}}
  \else
    \def\fitpurebptable#1{#1}
  \fi
  \fitpurebptable{%
  \begin{tabular}{lllccrr}
    \toprule
    Model & Context & Method & Topology & Global batch & Tok/s & Peak HBM (GiB) \\
    \midrule
    NemotronDiffusion 14B & 64K & Best baseline & DP8/TP1/CP2 & 8 & 20{,}092 & 89.5 \\
                  &     & Pure BP            & DP8/TP1/BP2 & 8 & 15{,}223 & 110.7 \\
                  &     & Ours               & DP8/TP1/CP2/BP2 & 8 & \textbf{22{,}152} & 89.8 \\
    \addlinespace
    NemotronDiffusion 14B & 128K & Best baseline & DP4/TP1/CP4 & 4 & 13{,}114 & 90.0 \\
                  &      & Pure BP            & DP4/TP1/BP4 & 4 & OOM & OOM \\
                  &      & Ours               & DP4/TP1/CP4/BP4 & 4 & \textbf{15{,}136} & 90.1 \\
    \midrule
    DiffusionGemma 26B-A4B & 64K & Best baseline & DP4/TP1/EP2/CP2 & 8 & 22{,}154 & 109.4 \\
                   &     & Pure BP            & DP4/TP1/EP2/BP2 & 8 & OOM & OOM \\
                   &     & Ours               & DP4/TP1/EP2/CP2/BP2 & 8 & \textbf{27{,}653} & 113.5 \\
    \addlinespace
    DiffusionGemma 26B-A4B & 128K & Best baseline & DP2/TP1/EP2/CP4 & 4 & 14{,}584 & 112.0 \\
                   &      & Pure BP            & DP2/TP1/EP2/BP4 & 4 & OOM & OOM \\
                   &      & Ours               & DP2/TP1/EP2/CP4/BP4 & 4 & \textbf{21{,}049} & 113.6 \\
    \bottomrule
  \end{tabular}}

\end{table}


\section{Load Balancing Analysis}
\label{app:dual-end-scheduling-ablation}

Load balancing delivers substantial gains without increasing memory use
(Table~\ref{tab:dual-end-scheduling-ablation}). Replacing contiguous block
assignment with early--late pairing improves throughput by $1.34\times$ for
NemotronDiffusion 3B and $1.28\times$ for DiffusionGemma 26B-A4B at 256K. Peak
HBM is unchanged for NemotronDiffusion and falls by 6.5\,GiB for DiffusionGemma.
Only the block assignment changes; the objective and
remaining training settings are fixed. These gains show that equal block
counts are not enough: balancing attention work across ranks is important
to realizing CSBP's throughput benefit.

% Two-node dual-end scheduling ablation collected on the Next Concept Modal
% account. DiffusionGemma uses the native all-block objective from clean commit
% e569364ae29ced06c4f457dfda5897aee120e488.
% NemotronDiffusion 3B, dual-end / contiguous:
% /persistent/runs/nemotron_labs_diffusion_3b_262144_block256_h200_nodes2_tp1_ep1_cp4_bp4_mb1_scheduledual_end_xnodecp_20260831_043616
% /persistent/runs/nemotron_labs_diffusion_3b_262144_block256_h200_nodes2_tp1_ep1_cp4_bp4_mb1_schedulecontiguous_xnodecp_20260831_044243
% DiffusionGemma 26B-A4B, dual-end / contiguous:
% /persistent/runs/diffusiongemma_26b_a4b_it_262144_block256_h200_nodes2_tp1_ep2_cp8_bp8_mb1_scheduledual_end_xnodecp_20260914_001756
% /persistent/runs/diffusiongemma_26b_a4b_it_262144_block256_h200_nodes2_tp1_ep2_cp8_bp8_mb1_schedulecontiguous_xnodecp_20260914_025123
\begin{table}[H]
  \centering
  \caption{Load balancing accelerates CSBP at 256K context.}
  \label{tab:dual-end-scheduling-ablation}
  \small
  \setlength{\tabcolsep}{4pt}
  \renewcommand{\arraystretch}{1.35}
  \resizebox{\linewidth}{!}{%
  \begin{tabular}{@{}llcrrrrr@{}}
    \toprule
    Model & Topology & \shortstack{Global\\batch} & ms/step & Tok/s & \shortstack{MFU\\(\%)} & \shortstack{Peak HBM\\(GiB)} & Speedup \\
    \midrule
    NemotronDiffusion 3B & DP4/TP1/CP4/BP4 & 4 &
    88{,}950.4 / \textbf{66{,}500.3} &
    11{,}788 / \textbf{15{,}768} & 27.8 / \textbf{37.2} &
    57.0 / 57.0 &
    \textbf{1.34$\times$} \\
    \midrule
    DiffusionGemma 26B-A4B & DP1/TP1/EP2/CP8/BP8 & 2 &
    43{,}722.0 / \textbf{34{,}134.6} &
    11{,}991 / \textbf{15{,}359} & 14.4 / \textbf{18.4} &
    120.1 / \textbf{113.6} & \textbf{1.28$\times$} \\
    \bottomrule
  \end{tabular}}
  \par\vspace{2pt}
  \parbox{\linewidth}{\footnotesize\raggedright Pairs: contiguous~/~dual-end. Speedup: dual-end~/~contiguous. Bold marks the better value.}
\end{table}


\section{AR-to-BDLM Conversion Results}
\label{app:fast-dllm-v2-results}

CSBP's benefits extend to converting autoregressive models into BDLMs
(Table~\ref{tab:fast-dllm-v2-context-scaling}). Under the Fast-dLLM v2
objective, both Qwen backbones train faster and use less peak HBM at every
evaluated context length. Speedup grows from $1.26$--$1.27\times$ at 64K to
$1.27$--$1.33\times$ at 256K, where CSBP also saves 9.6\,GiB of HBM.

% Qwen3.8-27B run provenance:
% 64K pair: /persistent/runs/qwen3.8_27b_65536_h200_fast_dllm_v2_nodes2_tp2_cp8_20260909_175016
% 128K pair: /persistent/runs/qwen3.8_27b_131072_h200_fast_dllm_v2_nodes2_tp2_cp8_20260909_175015
% 256K pair: /persistent/runs/qwen3.8_27b_262144_h200_fast_dllm_v2_nodes2_tp2_cp8_20260909_172222
% Qwen3.5-27B run provenance:
% 64K pair: /persistent/runs/qwen3.5_27b_65536_h200_fast_dllm_v2_nodes2_tp2_cp8_20260914_034354
% 128K pair: /persistent/runs/qwen3.5_27b_131072_h200_fast_dllm_v2_nodes2_tp2_cp8_20260914_035529
% 256K pair: /persistent/runs/qwen3.5_27b_262144_h200_fast_dllm_v2_nodes2_tp2_cp8_20260914_035528
\begin{table}[H]
  \centering
  \caption{CSBP accelerates AR-to-BDLM conversion from 64K to 256K context.}
  \label{tab:fast-dllm-v2-context-scaling}
  \small
  \setlength{\tabcolsep}{4pt}
  \renewcommand{\arraystretch}{1.35}
  \resizebox{\linewidth}{!}{%
  \begin{tabular}{@{}llllrrrr@{}}
    \toprule
    Model & Context & \shortstack{Best baseline\\topology} & \shortstack{Ours\\topology} & Tok/s & Speedup & \shortstack{MFU\\(\%)} & \shortstack{Peak HBM\\(GiB)} \\
    \midrule
    Qwen3.8-27B & 64K  & DP1/TP2/CP8 & DP1/TP2/CP8/BP8 & 4{,}697 / \textbf{5{,}905} & \textbf{1.26$\times$} & 22.3 / \textbf{28.0} & 64.2 / \textbf{62.5} \\
    Qwen3.8-27B & 128K & DP1/TP2/CP8 & DP1/TP2/CP8/BP8 & 3{,}742 / \textbf{4{,}853} & \textbf{1.30$\times$} & 21.4 / \textbf{27.8} & 83.0 / \textbf{78.3} \\
    Qwen3.8-27B & 256K & DP1/TP2/CP8 & DP1/TP2/CP8/BP8 & 2{,}623 / \textbf{3{,}490} & \textbf{1.33$\times$} & 20.1 / \textbf{26.8} & 120.7 / \textbf{111.1} \\
    \midrule
    Qwen3.5-27B & 64K  & DP1/TP2/CP8 & DP1/TP2/CP8/BP8 & 4{,}709 / \textbf{5{,}967} & \textbf{1.27$\times$} & 22.3 / \textbf{28.3} & 64.2 / \textbf{62.5} \\
    Qwen3.5-27B & 128K & DP1/TP2/CP8 & DP1/TP2/CP8/BP8 & 3{,}722 / \textbf{4{,}850} & \textbf{1.30$\times$} & 21.3 / \textbf{27.8} & 83.0 / \textbf{78.3} \\
    Qwen3.5-27B & 256K & DP1/TP2/CP8 & DP1/TP2/CP8/BP8 & 2{,}641 / \textbf{3{,}351} & \textbf{1.27$\times$} & 20.3 / \textbf{25.7} & 120.7 / \textbf{111.1} \\
    \bottomrule
  \end{tabular}}
  \par\vspace{2pt}
  \parbox{\linewidth}{\footnotesize\raggedright Pairs: best baseline~/~ours. Bold marks the better value.}
\end{table}


\section{Effectiveness Across Different Block Sizes}
\label{app:block-size-sensitivity}

CSBP maintains its advantage across an eightfold range of block sizes
(Table~\ref{tab:block-size-sensitivity}). From 128 to 1,024 tokens per block,
speedup stays within $1.14$--$1.15\times$ for NemotronDiffusion 14B and
$1.33$--$1.44\times$ for DiffusionGemma 26B-A4B.

% Each row is a matched best-baseline (pure CP) / ours (fused CP+BP)
% comparison. DiffusionGemma uses the native all-block objective from clean
% commit e569364ae29ced06c4f457dfda5897aee120e488. The block-256 pair remains
% the canonical fixed-topology 128K measurement; Table 1 reports 256K results.
% NemotronDiffusion 14B provenance (best baseline / ours):
% 128: /persistent/runs/nemotron_labs_diffusion_14b_131072_block128_h200_nodes2_tp1_ep1_cp4_bp4_mb1_20260830_195154
%      /persistent/runs/nemotron_labs_diffusion_14b_131072_block128_h200_nodes2_tp1_ep1_cp4_bp4_mb1_20260830_111859
% 256: /persistent/runs/nemotron_labs_diffusion_14b_131072_block256_h200_nodes2_tp1_ep1_cp4_bp4_mb1_xnodecp_20260828_185802
% 512: /persistent/runs/nemotron_labs_diffusion_14b_131072_block512_h200_nodes2_tp1_ep1_cp4_bp4_mb1_20260830_195154
%      /persistent/runs/nemotron_labs_diffusion_14b_131072_block512_h200_nodes2_tp1_ep1_cp4_bp4_mb1_20260830_111849
% 1024:/persistent/runs/nemotron_labs_diffusion_14b_131072_block1024_h200_nodes2_tp1_ep1_cp4_bp4_mb1_20260830_195240
%      /persistent/runs/nemotron_labs_diffusion_14b_131072_block1024_h200_nodes2_tp1_ep1_cp4_bp4_mb1_20260830_113002
% DiffusionGemma provenance (best baseline / ours):
% 128: /persistent/runs/diffusiongemma_26b_a4b_it_131072_block128_h200_nodes2_tp1_ep2_cp4_bp4_mb1_scheduledual_end_xnodecp_20260914_025534
% 256: /persistent/runs/diffusiongemma_26b_a4b_it_131072_block256_h200_nodes2_tp1_ep2_cp4_bp4_mb1_scheduledual_end_xnodecp_20260914_000552
% 512: /persistent/runs/diffusiongemma_26b_a4b_it_131072_block512_h200_nodes2_tp1_ep2_cp4_bp4_mb1_scheduledual_end_xnodecp_20260914_025535
% 1024:/persistent/runs/diffusiongemma_26b_a4b_it_131072_block1024_h200_nodes2_tp1_ep2_cp4_bp4_mb1_scheduledual_end_xnodecp_20260914_025535
\begin{table}[H]
  \centering
  \caption{CSBP sustains speedups across block sizes at 128K context.}
  \label{tab:block-size-sensitivity}
  \small
  \setlength{\tabcolsep}{5pt}
  \renewcommand{\arraystretch}{1.35}
  \resizebox{\linewidth}{!}{%
  \begin{tabular}{@{}lrllcrrrrr@{}}
    \toprule
    Model & \shortstack{Block\\size} & \shortstack{Best\\baseline\\topology} & \shortstack{Ours\\topology} & \shortstack{Global\\batch} & ms/step & Tok/s & \shortstack{MFU\\(\%)} & \shortstack{Peak HBM\\(GiB)} & Speedup \\
    \midrule
    NemotronDiffusion 14B & 128 & DP4/TP1/CP4 & DP4/TP1/CP4/BP4 & 4 &
    39{,}187.3 / \textbf{34{,}279.4} & 13{,}379 / \textbf{15{,}295} &
    34.3 / \textbf{39.2} & \textbf{90.0} / 90.1 & \textbf{1.14$\times$} \\
    NemotronDiffusion 14B & 256 & DP4/TP1/CP4 & DP4/TP1/CP4/BP4 & 4 &
    39{,}979.2 / \textbf{34{,}637.4} & 13{,}114 / \textbf{15{,}136} &
    33.7 / \textbf{38.9} & \textbf{90.0} / 90.1 & \textbf{1.15$\times$} \\
    NemotronDiffusion 14B & 512 & DP4/TP1/CP4 & DP4/TP1/CP4/BP4 & 4 &
    38{,}852.1 / \textbf{34{,}197.3} & 13{,}494 / \textbf{15{,}331} &
    34.7 / \textbf{39.4} & \textbf{90.0} / 90.1 & \textbf{1.14$\times$} \\
    NemotronDiffusion 14B & 1{,}024 & DP4/TP1/CP4 & DP4/TP1/CP4/BP4 & 4 &
    39{,}201.0 / \textbf{34{,}392.4} & 13{,}374 / \textbf{15{,}244} &
    34.5 / \textbf{39.3} & \textbf{90.0} / 90.1 & \textbf{1.14$\times$} \\
    \midrule
    DiffusionGemma 26B-A4B & 128 & DP2/TP1/EP2/CP4 & DP2/TP1/EP2/CP4/BP4 & 4 &
    33{,}271.3 / \textbf{24{,}844.4} & 15{,}758 / \textbf{21{,}103} &
    12.6 / \textbf{16.9} & \textbf{111.8} / 113.6 & \textbf{1.34$\times$} \\
    DiffusionGemma 26B-A4B & 256 & DP2/TP1/EP2/CP4 & DP2/TP1/EP2/CP4/BP4 & 4 &
    35{,}949.5 / \textbf{24{,}907.5} & 14{,}584 / \textbf{21{,}049} &
    11.7 / \textbf{16.9} & \textbf{112.0} / 113.6 & \textbf{1.44$\times$} \\
    DiffusionGemma 26B-A4B & 512 & DP2/TP1/EP2/CP4 & DP2/TP1/EP2/CP4/BP4 & 4 &
    33{,}319.8 / \textbf{24{,}829.4} & 15{,}735 / \textbf{21{,}116} &
    12.7 / \textbf{17.0} & \textbf{112.1} / 113.7 & \textbf{1.34$\times$} \\
    DiffusionGemma 26B-A4B & 1{,}024 & DP2/TP1/EP2/CP4 & DP2/TP1/EP2/CP4/BP4 & 4 &
    33{,}033.5 / \textbf{24{,}850.7} & 15{,}871 / \textbf{21{,}097} &
    12.9 / \textbf{17.1} & \textbf{112.3} / 113.7 & \textbf{1.33$\times$} \\
    \bottomrule
  \end{tabular}}
  \par\vspace{2pt}
  \parbox{\linewidth}{\footnotesize\raggedright Pairs: best baseline~/~ours. Bold marks the better value.}
\end{table}


\section{Savings Breakdowns}
\label{app:sources-of-speedup}

Table~\ref{tab:operator-profile-256k} separates attention communication from
end-to-end step time. NemotronDiffusion 14B uses full attention, so keeping
corrupted K/V local halves its logical attention traffic; overlap limits the
reduction in exposed attention communication to 20.4\%. DiffusionGemma 26B-A4B
has a second, architecture-specific source of savings: 25 of its 30 attention
layers use a 1,024-token sliding window. The tested CP8 baseline uses full-shard
K/V exchange in all 30 layers. CSBP keeps corrupted K/V local and, in the 25
sliding-window layers, its window-aware exchange sends only the clean rows
needed by local queries; only the five global-attention layers exchange
complete clean K/V shards. This reduces logical attention traffic by 93.5\%
and exposed attention communication by 93.4\%.

% End-to-end step time reuses the selected 256K throughput measurements in
% tables/main_standard_bdlm_256k.tex. Operator rows are per-GPU means from the
% bounded step-4 CUPTI traces of matched runs on 16 H200 GPUs:
% NemotronDiffusion 14B:
% /persistent/runs/nemotron_labs_diffusion_14b_262144_block256_h200_nodes2_tp1_ep1_cp8_bp8_mb1_xnodecp_20260831_061523
% DiffusionGemma 26B-A4B:
% CP: /persistent/runs/diffusiongemma_26b_a4b_it_262144_block256_h200_nodes2_tp1_ep2_cp8_bp8_mb1_scheduledual_end_xnodecp_20260914_030555
% CSBP: /persistent/runs/diffusiongemma_26b_a4b_it_262144_block256_h200_nodes2_tp1_ep2_cp8_bp8_mb1_scheduledual_end_xnodecp_20260914_030640
\begin{table}[H]
  \centering
  \caption{Attention and step-time breakdown at 256K: CP8 baseline vs. CSBP.}
  \label{tab:operator-profile-256k}
  \small
  \setlength{\tabcolsep}{3pt}
  \renewcommand{\arraystretch}{1.12}
  \begin{tabular}{lcc}
    \toprule
    Measure & \shortstack{NemotronDiffusion 14B\\CP8 $\rightarrow$ CP8/BP8} & \shortstack{DiffusionGemma 26B-A4B\\CP8 $\rightarrow$ CP8/BP8} \\
    \midrule
    Step time (s)
      & 66.50 $\rightarrow$ 56.37 \textcolor{green!50!black}{\textbf{($1.18\times$)}}
      & 49.66 $\rightarrow$ 34.13 \textcolor{green!50!black}{\textbf{($1.45\times$)}} \\
    Forward GPU span (s)
      & 17.57 $\rightarrow$ 13.00 \textcolor{green!50!black}{\textbf{($-26.0\%$)}}
      & 10.99 $\rightarrow$ 5.83 \textcolor{green!50!black}{\textbf{($-47.0\%$)}} \\
    Backward GPU span (s)
      & 49.46 $\rightarrow$ 45.02 \textcolor{green!50!black}{\textbf{($-9.0\%$)}}
      & 36.50 $\rightarrow$ 26.62 \textcolor{green!50!black}{\textbf{($-27.1\%$)}} \\
    Attention traffic (GiB/GPU)
      & 210.00 $\rightarrow$ 105.00 \textcolor{green!50!black}{\textbf{($-50.0\%$)}}
      & 433.13 $\rightarrow$ 28.30 \textcolor{green!50!black}{\textbf{($-93.5\%$)}} \\
    Exposed attention comm. (s/GPU)
      & 3.034 $\rightarrow$ 2.414 \textcolor{green!50!black}{\textbf{($-20.4\%$)}}
      & 10.939 $\rightarrow$ 0.723 \textcolor{green!50!black}{\textbf{($-93.4\%$)}} \\
    \bottomrule
  \end{tabular}
\end{table}


The 93\% reductions apply to attention communication, not the complete
training step. DiffusionGemma's step time falls by 31.3\%, while dense and
expert GEMM durations remain nearly unchanged. Both forward and backward
become faster, yielding $1.18\times$ and $1.45\times$ end-to-end speedup for
NemotronDiffusion 14B and DiffusionGemma 26B-A4B, respectively.

\section{Complete H200 Profiles}
\label{app:h200-profiles}

CSBP improves the throughput--memory trade-off across the tested models and
context lengths (Table~\ref{tab:h200-two-node-1p2}). The independently selected
configurations deliver speedups up to $1.61\times$. Even when throughput is
nearly tied, context sharding can substantially reduce memory: for NemotronDiffusion
3B at 128K, CSBP stays within 1\% of the DP16 baseline's throughput while
saving 41\,GiB of peak HBM. This case highlights another practical benefit
of CSBP: retaining throughput while freeing memory.

% Requires \usepackage{booktabs,graphicx}.
% The original candidate matrix and later throughput-neighbor profiles are
% documented with run identifiers at github.com/tarsur909/bdlm_parallel/issues/26.
% NemotronDiffusion 3B 128K best-baseline DP16 profile:
% /persistent/runs/nemotron_labs_diffusion_3b_131072_block256_h200_nodes2_tp1_ep1_cp1_bp1_mb1_20260829_080718.
% NemotronDiffusion 3B 128K ours profile:
% /persistent/runs/nemotron_labs_diffusion_3b_131072_block256_h200_nodes2_tp1_ep1_cp2_bp2_mb1_xnodecp_20260813_140602.
% NemotronDiffusion 3B 256K selected DP4/CP4 pair:
% /persistent/runs/nemotron_labs_diffusion_3b_262144_block256_h200_nodes2_tp1_ep1_cp4_bp4_mb1_xnodecp_20260830_201529.
% NemotronDiffusion 3B 512K ours/best-baseline pair:
% /persistent/runs/nemotron_labs_diffusion_3b_524288_block256_h200_nodes2_tp1_ep1_cp8_bp8_mb1_xnodecp_20260813_133650.
% DiffusionGemma 64K degree-2 pair:
% /persistent/runs/diffusiongemma_26b_a4b_it_65536_block256_h200_nodes2_tp1_ep2_cp2_bp2_mb1_scheduledual_end_xnodecp_20260913_233237.
% DiffusionGemma 128K degree-4 pair:
% /persistent/runs/diffusiongemma_26b_a4b_it_131072_block256_h200_nodes2_tp1_ep2_cp4_bp4_mb1_scheduledual_end_xnodecp_20260914_000552.
% DiffusionGemma 256K degree-8 pair:
% /persistent/runs/diffusiongemma_26b_a4b_it_262144_block256_h200_nodes2_tp1_ep2_cp8_bp8_mb1_scheduledual_end_xnodecp_20260914_001756.
% DiffusionGemma 512K degree-8 pair with TP2/SP:
% /persistent/runs/diffusiongemma_26b_a4b_it_524288_block256_h200_nodes2_tp2_ep1_cp8_bp8_mb1_scheduledual_end_xnodecp_20260914_012613.
% HBM is max_peak_mib / 1024, rounded to one decimal, where max_peak_mib is
% the maximum peak allocated memory across all 16 ranks.
\begin{table}[H]
  \centering
  \caption{Throughput and memory gains across long-context training configurations.}
  \label{tab:h200-two-node-1p2}
  \small
  \setlength{\tabcolsep}{5pt}
  \renewcommand{\arraystretch}{1.15}
  \resizebox{\linewidth}{!}{%
  \begin{tabular}{@{}llllcrrrrr@{}}
    \toprule
    Model & Context & \shortstack{Ours\\topology} & \shortstack{Best\\baseline\\topology} & \shortstack{Global\\batch} & Tok/s & \shortstack{MFU\\(\%)} & \shortstack{Peak HBM\\(GiB)} & \shortstack{HBM saved\\(GiB)} & Speedup \\
    \midrule
%    NemotronDiffusion 3B & 128K & DP2/TP1/CP8/BP8 & DP2/TP1/CP8 & 2 & 28{,}698 / 23{,}987 & 37.2 / 31.1 & 25.2 / 26.2 & 1.0 & \textbf{1.20$\times$} \\
%    NemotronDiffusion 3B & 128K & DP4/TP1/CP4/BP4 & DP4/TP1/CP4 & 4 & 28{,}730 / 24{,}136 & 37.3 / 31.3 & 35.6 / 35.7 & 0.1 & \textbf{1.19$\times$} \\
    NemotronDiffusion 3B & 128K & DP8/TP1/CP2/BP2 & DP16/TP1 & 8 / 16 & 28{,}891 / 29{,}086 & 37.5 / 37.7 & 56.5 / 97.5 & 41.0 & 0.99$\times$ \\
    NemotronDiffusion 3B & 256K & DP4/TP1/CP4/BP4 & DP4/TP1/CP4 & 4 & 16{,}027 / 13{,}434 & 37.8 / 31.7 & 57.0 / 57.1 & 0.1 & \textbf{1.19$\times$} \\
    NemotronDiffusion 3B & 512K & DP2/TP1/CP8/BP8 & DP2/TP1/CP8 & 2 & 8{,}352 / 6{,}895 & 37.4 / 30.8 & 58.1 / 61.7 & 3.6 & \textbf{1.21$\times$} \\
    \midrule
%    NemotronDiffusion 8B & 64K  & DP8/TP1/CP2/BP2 & DP8/TP1/CP2 & 8  & 29{,}603 / 26{,}521 & 37.5 / 33.6 & 64.1 / 63.9 & $-0.2$ & \textbf{1.12$\times$} \\
    NemotronDiffusion 8B & 64K  & DP8/TP1/CP2/BP2 & DP8/TP1/CP2 & 16 & 29{,}829 / 26{,}669 & 37.8 / 33.8 & 96.9 / 96.4 & $-0.5$ & \textbf{1.12$\times$} \\
    NemotronDiffusion 8B & 128K & DP8/TP1/CP2/BP2 & DP8/TP1/CP2 & 8  & 19{,}653 / 16{,}580 & 38.5 / 32.5 & 96.9 / 96.4 & $-0.5$ & \textbf{1.19$\times$} \\
    NemotronDiffusion 8B & 256K & DP4/TP1/CP4/BP4 & DP4/TP1/CP4 & 4  & 11{,}379 / 9{,}668 & 38.0 / 32.3 & 97.5 / 97.5 & 0.0 & \textbf{1.18$\times$} \\
    NemotronDiffusion 8B & 512K & DP2/TP1/CP8/BP8 & DP2/TP1/CP8 & 2  & 6{,}254 / 5{,}202 & 38.2 / 31.8 & 98.5 / 99.5 & 1.0 & \textbf{1.20$\times$} \\
    \midrule
    NemotronDiffusion 14B & 64K  & DP8/TP1/CP2/BP2 & DP8/TP1/CP2 & 8 & 22{,}152 / 20{,}092 & 38.8 / 35.2 & 89.8 / 89.5 & $-0.3$ & \textbf{1.10$\times$} \\
    NemotronDiffusion 14B & 128K & DP4/TP1/CP4/BP4 & DP4/TP1/CP4 & 4 & 15{,}136 / 13{,}114 & 38.9 / 33.7 & 90.1 / 90.0 & $-0.1$ & \textbf{1.15$\times$} \\
    NemotronDiffusion 14B & 256K & DP2/TP1/CP8/BP8 & DP2/TP1/CP8 & 2 & 9{,}301 / 7{,}884 & 39.0 / 33.1 & 90.6 / 91.1 & 0.5 & \textbf{1.18$\times$} \\
    NemotronDiffusion 14B & 512K & DP1/TP2/CP8/BP8 & DP1/TP2/CP8 & 1 & 5{,}052 / 4{,}242 & 37.7 / 31.6 & 71.9 / 72.7 & 0.8 & \textbf{1.19$\times$} \\
    \midrule
    DiffusionGemma 26B-A4B & 64K  & DP4/TP1/EP2/CP2/BP2 & DP4/TP1/EP2/CP2 & 8 & 27{,}653 / 22{,}154 & 16.0 / 12.8 & 113.5 / 109.4 & $-4.1$ & \textbf{1.25$\times$} \\
    DiffusionGemma 26B-A4B & 128K & DP2/TP1/EP2/CP4/BP4 & DP2/TP1/EP2/CP4 & 4 & 21{,}049 / 14{,}584 & 16.9 / 11.7 & 113.6 / 112.0 & $-1.6$ & \textbf{1.44$\times$} \\
    DiffusionGemma 26B-A4B & 256K & DP1/TP1/EP2/CP8/BP8 & DP1/TP1/EP2/CP8 & 2 & 15{,}359 / 10{,}557 & 18.4 / 12.7 & 113.6 / 118.1 & 4.6 & \textbf{1.45$\times$} \\
    DiffusionGemma 26B-A4B & 512K & DP1/TP2/CP8/BP8/SP & DP1/TP2/CP8/SP & 1 & 9{,}815 / 6{,}079 & 20.3 / 12.6 & 106.5 / 127.4 & 20.9 & \textbf{1.61$\times$} \\
    \bottomrule
  \end{tabular}}
  \par\vspace{2pt}
  \parbox{\linewidth}{\footnotesize\raggedright Pairs: ours~/~best baseline. HBM saved: best baseline minus ours.}
\end{table}

